\section{System Model}
\label{sec:model}

A guest bursts batch GPU pods onto a shared cluster under a fair-use policy that
allows at most $K$ concurrent pods and requires each admitted pod to sustain a
utilization floor. Time is slotted into control epochs $t\in\mathbb{N}$. Let
$a_t\in\{0,\dots,C\}$ be the capacity the cluster can productively give the guest
at epoch $t$. It is set by the other tenants, so it is exogenous, stochastic and
not directly observable. The guest chooses a window $w_t\in\{0,\dots,K\}$. Useful
goodput is $G_t=\min(w_t,a_t)$ and the surplus $(w_t-a_t)^{+}$ is over-admitted.
We are deliberately careful about what that surplus costs, because the obvious
reading is wrong. A pod that requests a GPU when none is allocatable does not run
at low utilization, it stays pending, so surplus admissions are not per-pod
utilization violations. Their cost is that a pending claim captures capacity the
moment the neighbour releases it, and that a guest holding many such claims is
what fair-use review looks at. We therefore treat $\rho$ as a measure of
\emph{claims in excess of available capacity} and validate it against the quantity
it is supposed to proxy, namely the throughput the co-tenant keeps
(\S\ref{sec:eval-contended}). The tolerance $\epsilon$ is a guest-chosen budget on
that excess rather than a figure read off the published policy.
Writing $\rho\triangleq\limsup_T\frac1T\sum_{t<T}\mathbf 1[w_t>a_t]$ for the
long-run over-admission fraction, the guest solves
\begin{equation}\label{eq:polite}
\max_{\{w_t\}}\ \mathbb{E}[\min(w_t,a_t)]
\quad\text{s.t.}\quad
\underbrace{w_t\le K\ \forall t}_{\text{hard cap}},\qquad
\underbrace{\rho\le\epsilon}_{\text{policy}} .
\end{equation}
The constraint is two-sided, since the guest must stay below the cap and above the
floor, which already separates the problem from loss-only congestion control.

\noindent\textbf{The information-to-effect lag.}
Two distinct delays act on this loop and only their sum matters. A decision taken
at wall-clock time $\tau$ rests on a capacity reading of age
$D_{\mathrm{obs}}$, and it takes effect at $\tau+D_{\mathrm{act}}$, when the
admitted pod actually occupies the resource and its productivity is scored. The
state that determines the payoff is therefore the capacity at
$\tau+D_{\mathrm{act}}$, while the information the policy conditions on is the
capacity at $\tau-D_{\mathrm{obs}}$. Writing
\begin{equation}\label{eq:Ddef}
D\ \triangleq\ D_{\mathrm{obs}}+D_{\mathrm{act}}
\end{equation}
for that information-to-effect lag and relabelling $t=\tau+D_{\mathrm{act}}$, the
admissible policies are exactly the causal functions of $\{a_{s-D}\}_{s\le t}$.
This is the only sense in which $D$ is used below, and it is the sense the measurements
report. It is not the age of the record at the decision instant, which is only
$D_{\mathrm{obs}}$.

The distinction matters because the two components differ by two orders of
magnitude. Polling occupancy of resources already running costs
$D_{\mathrm{obs}}=27$--$33$\,ms, well under one epoch, so a paper that identified
$D$ with observation age alone would predict that feedback is always useful here.
Actuation is what dominates. A newly admitted pod needs
$D_{\mathrm{act}}\approx2.46$\,s of CUDA start-up before it occupies the GPU it was
admitted for (\S\ref{sec:eval}), and the neighbour keeps moving during that
interval. So $D$ is $2$ to $3$ epochs, not a fraction of one, and it is set by
actuation rather than by monitoring design.

That also answers the natural objection that $D$ is an engineering choice.
Distributed schedulers routinely hide actuation latency by launching more work
than they need and binding to whichever reservation lands first, discarding the
rest~\cite{sparrow}. Under a one-sided constraint that works. Under the two-sided
constraint of \eqref{eq:polite} it does not.

\begin{remark}[Why the actuation term is hard to remove here]
\label{prop:irreducible}
For the controller studied here, where a scale-out decision creates a new CUDA
worker, the measured information-to-effect lag is about $2.46$\,s
(\S\ref{sec:eval-detection}) and is dominated by worker start-up. The usual way to
hide that latency is to keep spare admitted capacity, whether by speculative
replicas that are later cancelled or by an idle standby, and both consume claims
against the guest's own window while contributing nothing, which is what the
guest is trying to avoid. We do \emph{not} claim this is irreducible in general.
A persistent worker pulling from a queue, pipelined replacement that launches a
successor while its predecessor still works, or a warm pool all reduce or remove
the start-up term, at the cost of holding capacity. Which trade is preferable is a
design question, and the results below take $D$ as measured for this design rather
than as a universal floor.
\end{remark}
What separates this setting from load balancing on stale state, where the classical
prescription is to sample few and avoid acting greedily on old
readings~\cite{mitzenmacher,goel}, is that there a guest which does nothing is
compliant, so holding spare capacity to hide latency is free. Here it is not.
Everything that follows treats $D$ as measured and asks what is achievable at that
lag.

\noindent\textbf{The contended unit.}
Politeness is decided at the margin, so we fix attention on the single contended
resource unit whose availability fluctuates with the neighbour. Let
$x_t\in\{0,1\}$ indicate that this unit is free ($x_t=1$) or taken ($x_t=0$), so
admitting it when taken is the over-admission the policy punishes. We take $x_t$
stationary with $\Pr[x_t=1]=\tfrac12$, a fully contended unit, and note that the
asymmetric case only rescales the constants below. The guest's marginal decision
$m_t\in\{0,1\}$ is a possibly randomized causal function of the age-$D$
observation $\hat x_t\triangleq x_{t-D}$. Marginal goodput and over-admission are
\begin{equation}\label{eq:margin}
g \triangleq \Pr[m_t=1,\,x_t=1],\qquad
\rho \triangleq \Pr[m_t=1,\,x_t=0].
\end{equation}
Everything below concerns the achievable $(\rho,g)$ region under lag $D$. The
stationary $\rho$ of \eqref{eq:margin} is the ergodic limit of the epoch fraction
in \eqref{eq:polite}, so we reuse the symbol.

\begin{proposition}[Hard-cap safety]\label{prop:safety}
Any window update that clamps additive increase at $K$ and never increases $w_t$
on a congestion signal keeps $w_t\le K$ for every $t$ deterministically, so the
pod cap is never violated.
\end{proposition}
\begin{proof}
Induction on $t$. Increase is clamped to $K$ and decrease cannot exceed the
current value.
\end{proof}

\section{How Much Stale Feedback Is Worth}
\label{sec:limit}

We now characterize exactly how much closed-loop admission gains over a fixed
budget when the only signal is a $D$-old observation. The answer is one
interpretable quantity, the autocorrelation of capacity at the sensing lag.

\begin{lemma}[Predictive disagreement]\label{lem:q}
Let $q(D)\triangleq\Pr[x_t\neq x_{t-D}]$ be the lag-$D$ disagreement probability
and $r(D)\triangleq\mathbb{E}[(2x_t\!-\!1)(2x_{t-D}\!-\!1)]$ the lag-$D$
autocorrelation of centered availability. Then $r(D)=1-2q(D)\in[-1,1]$, and by
stationarity and symmetry the joint law of $(\hat x_t,x_t)$ is
$\Pr[\hat x_t{=}i,x_t{=}j]=\tfrac14\big(1+(2i{-}1)(2j{-}1)\,r(D)\big)$.
For a symmetric two-state Markov capacity with per-epoch flip probability
$p\in(0,\tfrac12)$ we have $r(D)=(1-2p)^{D}=e^{-D/\tau_c}$ with coherence time
$\tau_c\triangleq-1/\ln(1-2p)\approx 1/(2p)$.
\end{lemma}
\begin{proof}
$r=\Pr[\text{agree}]-\Pr[\text{disagree}]=1-2q$. The joint law follows from the
two symmetric marginals and the single correlation. For the Markov chain, writing
$\lambda_2\triangleq1-2p$ for the second eigenvalue, the $D$-step transition is
$P^D=\tfrac12\!\left(\begin{smallmatrix}1+\lambda_2^D&1-\lambda_2^D\\1-\lambda_2^D&1+\lambda_2^D\end{smallmatrix}\right)$,
so $\Pr[\text{disagree}]=\tfrac12(1-\lambda_2^D)$ and $r=\lambda_2^D$. The
two-state chain's autocorrelation and its total-variation distance to
stationarity are standard mixing facts~\cite{lpw}.
\end{proof}

\begin{theorem}[Achievable region and the feedback advantage]\label{thm:limit}
For the symmetric two-state Markov capacity of Lemma~\ref{lem:q}, under age-$D$
sensing the achievable $(\rho,g)$ region of \eqref{eq:margin} is the
quadrilateral with vertices
$\{(0,0),\,(\tfrac{q}{2},\tfrac{1-q}{2}),\,(\tfrac12,\tfrac12),\,
(\tfrac{1-q}{2},\tfrac{q}{2})\}$ with $q=q(D)$, whose efficient upper boundary is
the polyline $(0,0)\!\to\!(\tfrac q2,\tfrac{1-q}2)\!\to\!(\tfrac12,\tfrac12)$.
The best observation-oblivious policy is confined to the diagonal $g=\rho$.
Consequently the goodput gain of feedback over a fixed budget, maximized over
matched over-admission levels, is
\begin{equation}\label{eq:adv}
\boxed{\;\Delta(D)\;=\;\tfrac12\,r(D)\;=\;\tfrac12\big(1-2q(D)\big)\;}
\end{equation}
attained at the kink $\rho=q/2$ by the threshold policy $m_t=\hat x_t$, while at a
lower level $\rho<q/2$ the gain is the smaller $\rho(1-2q)/q$. The peak value of
closed-loop admission is therefore exactly half the autocorrelation of available
capacity at the sensing lag (Fig.~\ref{fig:theory}). Optimality of memoryless
threshold policies for symmetric two-state Markov observation is classical in
opportunistic access~\cite{myopicaccess}, and what \eqref{eq:adv} adds is the
exact value of the resulting advantage over the best fixed budget.
\end{theorem}
\begin{proof}
For the symmetric two-state Markov chain the freshest sample $\hat x_t=x_{t-D}$ is
a sufficient statistic for $x_t$ given the whole delayed history
$\{x_{s-D}\}_{s\le t}$, since the Markov property makes older samples
conditionally uninformative. An optimal causal age-$D$ policy therefore depends on
the history only through $\hat x_t$, so we may restrict without loss to policies
parameterized by a pair $(\theta_1,\theta_0)$ with
$\theta_i=\Pr[m_t{=}1\mid\hat x_t{=}i]$. Using Lemma~\ref{lem:q},
$g=\tfrac12\big(\theta_1(1{-}q)+\theta_0 q\big)$ and
$\rho=\tfrac12\big(\theta_1 q+\theta_0(1{-}q)\big)$. This linear map sends the
square $[0,1]^2$ to the quadrilateral with the stated vertices, whose upper
boundary is the stated hull. Maximizing $g-\lambda\rho$ turns on $\theta_1$ before
$\theta_0$ because $q<\tfrac12$ makes the free-observation state strictly more
goodput-efficient, so the maximum-goodput point for $\rho\le q/2$ is the kink
$(\tfrac q2,\tfrac{1-q}2)$, achieved by $m_t=\hat x_t$. A fixed policy ignores
$\hat x_t$, so $\theta_1=\theta_0=\theta$ and $g=\rho=\theta/2$, the diagonal. The
vertical distance at $\rho=q/2$ is
$\tfrac{1-q}{2}-\tfrac q2=\tfrac{1-2q}{2}=\tfrac12 r(D)$.
\end{proof}

\begin{figure}[t]
\centering
\resizebox{0.92\columnwidth}{!}{%
\begin{tikzpicture}[>=Stealth]
  % value units * 10 = cm.  q=q(D); drawn at q=0.2 (r=0.6).
  \draw[->] (0,0) -- (6.0,0) node[right] {$\rho$ (over-admission)};
  \draw[->] (0,0) -- (0,6.0) node[above] {$g$ (goodput)};
  \fill[black!8] (0,0) -- (1,4) -- (5,5) -- (4,1) -- cycle;
  \node[black!55] at (4.05,1.55) {\footnotesize achievable};
  \draw[very thick] (0,0) -- (1,4) -- (5,5);
  \draw[dashed] (0,0) -- (5.4,5.4);
  \node[sloped,above,black!70] at (3.25,3.25) {\footnotesize best fixed budget, $g=\rho$};
  \draw[dotted] (1,0) -- (1,4);
  \node[below] at (1,0) {\footnotesize $\tfrac{q}{2}$};
  \fill (1,4) circle (2pt);
  \node[above left=-1pt] at (1,4) {\footnotesize $\bigl(\tfrac{q}{2},\tfrac{1-q}{2}\bigr)$};
  \node[align=center,anchor=south west] at (1.55,4.45)
        {\footnotesize closed-loop optimum\\[-2pt]\footnotesize $m_t=\hat x_t$};
  \fill (1,1) circle (1.6pt);
  \draw[<->,thick] (1,1.16) -- (1,3.84);
  \node[right=2pt] at (1,2.5) {$\Delta(D)=\tfrac12 r(D)$};
  \fill (5,5) circle (1.6pt);
  \node[below right=1pt] at (5,5) {\footnotesize $\bigl(\tfrac12,\tfrac12\bigr)$};
\end{tikzpicture}}
\caption{Under lag-$D$ sensing the achievable $(\rho,g)$ pairs form the shaded
hull, drawn at $q=0.2$. Fixed budgets lie on the diagonal $g=\rho$, and the
closed-loop optimum sits a vertical $\Delta(D)=\tfrac12 r(D)$ above it
(Thm.~\ref{thm:limit}). The advantage is a distance, which is why it admits a
closed form.}
\label{fig:theory}
\end{figure}

\begin{corollary}[The advantage decays with the lag]\label{cor:imp}
$\Delta(D)\!\to\!0$ as $r(D)\!\to\!0$. For the symmetric two-state Markov model
$\Delta(D)=\tfrac12 e^{-D/\tau_c}$, so any causal controller at lag $D$ gains at
most $\tfrac12 e^{-D/\tau_c}$ over a fixed budget, which is below any tolerance
$\eta$ once $D\ge\tau_c\ln(1/2\eta)$. Closed-loop admission helps only to
the extent that the guest senses the neighbour faster than the neighbour changes.
In this model low autocorrelation and unpredictability coincide, but in general
they need not, since a decorrelated process with $r(D){=}0$ can still be exactly
predictable from its history, as a random-phase deterministic square wave is. For
an arbitrary capacity process the vanishing of the advantage is therefore governed
by the $\beta$-mixing coefficient rather than by $r(D)$. At $D=0$ we have
$\Delta=\tfrac12$, and the advantage halves for every $\tau_c\ln 2$ epochs of
added age (Fig.~\ref{fig:theory}). App.~\ref{app:beta} makes the general
statement precise and extends the impossibility to arbitrary non-binary and
non-Markov capacity processes.
\end{corollary}

\noindent\textbf{Provenance, stated plainly.}
Equation \eqref{eq:adv} is elementary and we would rather derive it twice than
appear to have hidden how short it is. For a symmetric binary hypothesis the Bayes
error is $\tfrac12(1-\|\cdot\|_{\mathrm{TV}})$, so the largest advantage any test
holds over blind guessing is exactly half the total variation between the two
conditionals~\cite{tsybakov}. For a two-state chain that total variation is
$|\lambda_2|^{D}$, which is $r(D)$. Compose the two and \eqref{eq:adv} follows in
two lines. Each ingredient is old. Value of information is nonnegative and convex
in the prior~\cite{howard1966} and is bounded by payoff oscillation times the
movement of beliefs~\cite{delaragossner}. That movement at age $D$ is the
$\beta$-mixing coefficient~\cite{bradley2005}, and the mixing facts for a
two-state chain are standard~\cite{lpw}. Threshold optimality for positively
correlated two-state availability is Ahmad et al.~\cite{myopicaccess}, and
threshold optimality \emph{under state-information lag} is Kim~\cite{kim1985}.
Measuring the value of information as degradation from the perfect-information
gain is Brooks and Leondes~\cite{brooks1972}. Reducing a delay-$D$ decision
problem to an undelayed one by state augmentation is Katsikopoulos and
Engelbrecht~\cite{katsikopoulos}.

So the optimal policy is not ours and neither is the formula. Three things are.
First, the statement is a \emph{converse over every causal controller} at
\emph{matched over-admission}, not the value of one policy. The matched-constraint
normalization is what makes ``advantage over the best fixed budget'' well posed at
all, and it is the actual technical content of Thm.~\ref{thm:limit}. Second,
App.~\ref{app:beta} lifts the converse off the binary Markov case to any
stationary capacity process. Third, and most important for a systems venue,
Rem.~\ref{prop:irreducible} records that the lag in this design is dominated by a
term the guest cannot remove without holding spare capacity, which is what makes a
lag-indexed limit the useful object here rather than a modelling convenience. The exogenous structure also matters. Unlike the
driven-plant setting of networked control, the admission action does not affect
capacity, and that is what collapses the value of feedback onto a single
predictability coefficient.

\begin{corollary}[The two contention profiles are two points on \eqref{eq:adv}]
\label{cor:exp}
The contention study of \S\ref{sec:eval} instantiates this two-state $C{=}2$
model. Structured square load has $\tau_c\gg D$, so $r(D)$ is near $1$ and
feedback is near-maximally useful. Memoryless Poisson load drives
$\tau_c\lesssim D$, so $r(D)$ is near $0$ and feedback collapses onto the fixed
budget, which is the pre-registered null. The single curve \eqref{eq:adv} predicts
both cases and the crossover between them.
\end{corollary}

\section{An Adaptive Controller}
\label{sec:controller}

Theorem~\ref{thm:limit} leaves a shortfall that no fixed controller closes.
Additive-increase and multiplicative-decrease tracks a slowly varying neighbour
and approaches the kink when $\tau_c\gg D$, but over-admits when
$\tau_c\lesssim D$, since it keeps probing a signal that has become noise. A fixed
budget is safe everywhere and forfeits $\Delta(D)$ when $\tau_c\gg D$. Which of
the two is right depends on a quantity the controller can measure, because $r(D)$
is an autocorrelation of the probe stream the controller already collects.

Algorithm~\ref{alg:adaptive} estimates it online and switches. Beyond the AIMD
window it holds a flip count, a sample count, the previous reading and a $D$-deep
ring of past windows, so its state is $O(D)$ and its per-epoch cost is constant.
The window $w$ is a real-valued target and concurrency is $\lfloor w\rfloor$, which
matters because $\alpha<1$ otherwise leaves the window unable to move. It inherits
Prop.~\ref{prop:safety} and therefore never exceeds the pod cap. The counters
accumulate over the whole run rather than over a sliding window. That matters for
Lemma~\ref{thm:adapt}, since a fixed window of length $n$ leaves a floor on the
misclassification probability and the shortfall would then not vanish. A sliding
window is the right choice under drift and is what \S\ref{sec:eval} uses, at the
cost of the asymptotic statement.

\begin{algorithm}[t]
\caption{Adaptive admission for one contended unit}
\label{alg:adaptive}
\begin{algorithmic}[1]
\Require cap $K$, floor $w_{\min}$, step $\alpha$, backoff $\beta$, switch level
$\theta$, lag $D$, fixed budget $\bar a$
\State $w \gets w_{\min}$;\ $F \gets 0$;\ $n \gets 0$;\ $\hat a_{\mathrm{prev}} \gets \bot$
\For{each control epoch $t$}
  \State $\hat a \gets$ newest capacity reading \Comment{age $D_{\mathrm{obs}}$; a decision now lands at lag $D$}
  \If{$\hat a_{\mathrm{prev}} \neq \bot$}
    \State $F \gets F + \mathbf 1[\hat a \neq \hat a_{\mathrm{prev}}]$;\ \ $n \gets n+1$
  \EndIf
  \State $\hat a_{\mathrm{prev}} \gets \hat a$
  \State $\hat p \gets \min\!\big(F/\max(n,1),\ \tfrac12\big)$
    \Comment{clamp, else $(1-2\hat p)^{D}$ turns positive again for even $D$}
  \State $\hat r \gets (1-2\hat p)^{D}$
  \If{$\hat r > \theta$} \Comment{trackable, so run the loop}
    \State $w_{\mathrm{eff}} \gets$ window admitted $D$ epochs ago
      \Comment{from a $D$-deep ring of past windows}
    \If{$w_{\mathrm{eff}} > \hat a$} \Comment{that admission over-admitted}
      \State $w \gets \max(w_{\min},\ \lceil \beta w \rceil)$
    \Else
      \State $w \gets \min(w + \alpha,\ K)$
    \EndIf
  \Else \Comment{not trackable, so stop paying for the loop}
    \State $w \gets \min(\lfloor \bar a \rfloor,\ K)$
  \EndIf
  \State push $w$ onto the ring;\ \ admit $\lfloor w \rfloor$ pods
    \Comment{$w$ is a real-valued target, concurrency is its floor}
\EndFor
\end{algorithmic}
\end{algorithm}

The tracking analysis of Algorithm~\ref{alg:adaptive} assumes \textbf{(A1)}
stationary capacity $a_t\equiv a$ with $w_{\min}\le a\le K$, \textbf{(A2)}
reliable detection of an overshoot within $D$ epochs, and \textbf{(A3)} fixed
$\alpha,\beta,w_{\min},K$.

\begin{proposition}[Sawtooth realization of the kink]\label{prop:bound}
Under stationary capacity $a$, assumptions (A1) to (A3), a backoff deep enough to
clear the overshoot so that $\beta(a+D\alpha)\le a$, and ignoring the integer
rounding of the window, the trackable branch of Algorithm~\ref{alg:adaptive} forms
a periodic sawtooth about $a$ with long-run over-admission $\rho=D\alpha/\big((1-\beta)(a+D\alpha)\big)\approx
D\alpha/((1-\beta)a)$. Choosing $\alpha\le\epsilon(1-\beta)a/D$ gives
$\rho\le\epsilon$, and $\rho\to0$ as $D\to0$.
\end{proposition}
\begin{proof}[Proof sketch]
After a backoff the window restarts at $\beta(a+D\alpha)$ and rises by $\alpha$
per clean epoch, overshooting $a$ for $D$ epochs before detection triggers the
next backoff at peak $a+D\alpha$. The fraction of violating epochs per cycle gives $\rho$. Rounding the window to
integers perturbs the cycle length by at most one epoch and the bound by
$O(\alpha/a)$. This is the dynamic approach to the static-capacity kink of
Thm.~\ref{thm:limit}, and $D$ enters both here and in \eqref{eq:adv}, which is the
hinge between the systems measurement and the theory.
\end{proof}

Learning which branch to take has a sharp rate, and it is not the rate a generic
chain concentration bound suggests. For the symmetric two-state chain the flip
indicators $\mathbf 1[a_t\neq a_{t-1}]$ are \emph{exactly} independent
Bernoulli$(p)$ variables, since the flip probability is $p$ from either state and
after every history. Estimating $p$ therefore carries no mixing penalty, and the
relevant variance is $p(1-p)$ rather than the $1/\gamma$ proxy that a
reversible-chain Hoeffding bound inserts. Since $\gamma=2\min(p,1-p)$, that true
variance is about $\gamma/2$, so the spectral gap belongs in the numerator. We
therefore use the Bernstein form of \cite{paulin2015}, which is stated in the
asymptotic variance and is tight here.

Write $p_\theta$ for the solution of $(1-2p)^D=\theta$, let
$s_\theta\triangleq 2D\,\theta^{(D-1)/D}$ be the sensitivity
$|\mathrm{d}r/\mathrm{d}p|$ at that point, and let
$\sigma_\theta^2\triangleq p_\theta(1-p_\theta)$. Let $L>0$ be the local slope in
$r$ of the difference between the two branches' values at $\theta$, which is
positive because Prop.~\ref{prop:bound} and Cor.~\ref{cor:imp} cross transversally
there.

What can be said sharply is how hard the \emph{decision} is, separately from what
either branch is worth. Choosing a branch means deciding on which side of $\theta$
the quantity $r(D)$ lies, and that is a statistical problem with a rate. We state
it as such, and are explicit below about the step we do not take.

\begin{lemma}[Branch-decision rate]\label{thm:adapt}
Consider deciding, from $N$ consecutive age-$D$ readings of a symmetric two-state
capacity process near $p_\theta$, whether $r(D)>\theta$. The local minimax error in
$r$ of any causal rule is
\begin{equation}\label{eq:rate}
\Theta\!\left(\frac{s_\theta\,\sigma_\theta}{\sqrt N}\right),
\qquad
s_\theta\sigma_\theta=D\,\theta^{\frac{D-1}{D}}\sqrt{1-\theta^{2/D}},
\end{equation}
and the estimator inside Algorithm~\ref{alg:adaptive} attains it. Maximizing over
the switch level,
\begin{equation}\label{eq:sup}
\max_{\theta\in(0,1)} s_\theta\sigma_\theta
=\Big(1-\tfrac1D\Big)^{\frac{D-1}{2}}\sqrt D\ \longrightarrow\ e^{-1/2}\sqrt D,
\end{equation}
attained at $\theta=\big((D-1)/D\big)^{D/2}$, which increases to $e^{-1/2}$. When
$r(D)$ is bounded away from $\theta$ the error decays as $O(1/N)$ provided the
counters accumulate over the run.
\end{lemma}
\begin{proof}[Proof sketch]
Upper bound. Bernstein for the independent flip indicators \cite{paulin2015} gives
$|\hat p_n-p|=O(\sigma_\theta/\sqrt n)$ with high probability, and the delta method
through $p\mapsto(1-2p)^D$, whose derivative at $p_\theta$ is $s_\theta$, gives
$|\hat r_n-r(D)|=O(s_\theta\sigma_\theta/\sqrt n)$. Lower bound. Take two processes
whose flip probabilities differ by $\delta$. Their per-sample Kullback-Leibler
divergence is $\delta^2/(2\sigma_\theta^2)$, so no rule separates them over $N$
samples unless $\delta\gtrsim\sigma_\theta/\sqrt N$, which through the same map
means correlations closer than $s_\theta\sigma_\theta/\sqrt N$ cannot be
distinguished. Le Cam's two-point argument \cite{tsybakov} gives \eqref{eq:rate},
and the two processes differ in flip probability by $O(N^{-1/2})$ so their spectral
gaps agree to that order. For \eqref{eq:sup}, substituting $u=1-2p$ gives
$s_\theta\sigma_\theta=D\,u^{D-1}\sqrt{1-u^2}$, maximized at $u^2=(D-1)/D$ where
$\sqrt{1-u^2}=1/\sqrt D$, and $\theta=u^D$.
\end{proof}

\noindent\textbf{What this gives, and what it does not.}
Lemma~\ref{thm:adapt} concerns the decision, not goodput. Converting it into a
goodput shortfall needs the value of each branch as a function of $p$, and we do
not derive those. Prop.~\ref{prop:bound} gives the AIMD sawtooth under
\emph{constant} capacity rather than an AIMD value curve under a Markov neighbour,
so we have no proof that the two branch values cross transversally at any
particular $\theta$, and we therefore claim no minimax optimality for
Algorithm~\ref{alg:adaptive} itself. What the lemma does fix is the shape of the
learning cost. The spectral gap cancels, since the flip indicators are independent
and carry no mixing penalty, so a rate written as $1/\sqrt{\gamma_0N}$ would have
the dependence inverted. The lag enters only as $\sqrt D$, far milder than the
exponential decay of the ceiling in \eqref{eq:adv}. And the hardest level to sit at
is intermediate rather than extreme, which is the qualitative reason a controller
gains from measuring $r(D)$ rather than committing to one branch in advance.
Deriving the branch value curves, and with them a performance theorem for
Algorithm~\ref{alg:adaptive}, is left open.

Equation \eqref{eq:rate} is exact as written. The convenient reading
$s_\theta\sigma_\theta\approx\theta\sqrt{2D\log(1/\theta)}$ holds only for large
$D$ at fixed $\theta$ and fails at small $D$, since at $D=1$ with $\theta\to0$ the
left side tends to $1$ while that expression tends to $0$.



\noindent\textbf{Scope.}
Three assumptions earn the closed form. First, the contended unit is binary and
symmetric. For any symmetric stationary $x_t$ the value $\tfrac12 r(D)$ is
achievable by the memoryless threshold $m_t=\hat x_t$, and it is the exact
advantage when $x_t$ is Markov, since then $\hat x_t$ is a sufficient statistic.
For a general non-Markov symmetric process a history-using controller can do
strictly better, as a random-phase square wave with $r(D){=}0$ stays exactly
predictable from its phase, so the equality is Markov-specific and the general
converse is the mixing bound of App.~\ref{app:beta}. Second, stationarity, since
$r(D)$ is a lag-$D$ autocorrelation. Under slow drift it is replaced by a local
estimate, which is what Algorithm~\ref{alg:adaptive} tracks. Third, fixed $D$. If
the freshest sample has a random age $A$ the controller can read, the advantage is
the age-averaged $\tfrac12\,\mathbb{E}[r(A)]$ and the impossibility persists
whenever $A\gg\tau_c$ almost surely.

Integer capacity beyond a single contended unit is left open. A layer-cake
decomposition of $\min(w,a)$ and $(w-a)^{+}$ suggests one binary decision per
unit, but the per-layer decisions must be kept nested to correspond to a scalar
window, the layerwise advantage is only an upper bound unless a monotonicity
condition holds, and scoring surplus by amount rather than by violation
probability changes the problem. We prefer to state the single-unit case exactly
than the general case loosely.
